% Part: first-order-logic
% Chapter: introduction
% Section: models-theories

\documentclass[../../../include/open-logic-section]{subfiles}

\begin{document}

\olfileid{fol}{int}{mod}

\olsection{مدلونه او تيورۍ}

کله چې د لومړۍ درجې منطق نحو او معنٰی پوهنه تعريف کړو، د
!!{structure}s او معنايي مفاهيمو د خاصيتونو څېړل پيلولے شو۔ د
!!{derivation} نظامونه هم تعريف او وڅېړلے شو۔ د !!{sentence}s د
يوه سټ په باب پوښتلے شو: کومې !!{structure}s د هغه سټ ټولې
!!{sentence}s رښتونې کوي؟ د !!{sentence}s يو سټ~$\Gamma$ چې راکړل
شي، هغه !!a{structure}~$\Struct{M}$ چې هغوئ پوره کوي د~$\Gamma$
\emph{مدل} بلل کېږي۔ ښايي له~$\Gamma$ څخه پيل او د هغې مدلونه
ولټوو---هغوئ څنګه دي؟ بايد څومره لوے يا واړه وي؟ خو ښايي له يوې
!!{structure} يا د !!{structure}s له يوې ټولګې څخه هم پيل او
پوښتنه وکړو: کومې !!{sentence}s پکښې رښتونې دي؟ ايا داسې
!!{sentence}s شته چې دا !!{structure}s په دې معنٰی
\emph{ځانګړې} کړي چې هغوئ، او يوازې هغوئ، پکښې رښتونې وي؟ دا ډول
پوښتنې د \emph{مدل تيورۍ} ډګر دے۔ همدغه پوښتنې د \emph{بديهي
طريقې} بنسټ هم جوړوي: د !!{structure}s د يوې ټولګې بيانول د
!!{sentence}s په يوه سټ، يعنې د يوې تيورۍ په بديهي اصولو۔ دا د دې
مشاهدې له امله ممکنېږي چې کټ مټ هغه !!{sentence}s د بديهي اصولو په
ټولو مدلونو کښې رښتونې دي چې د لومړۍ درجې منطق کښې له هغو بديهي
اصولو څخه لازمېږي۔

د يوې ډېرې ساده بېلګې په توګه مخکښ‌مرتبې وګورئ۔ مخکښ‌مرتبه پر کوم
سټ~$A$ داسې اړيکه~$R$ ده چې انعکاسي او انتقالي دواړه وي۔ سټ~$A$
چې دوه‌ځاييزه اړيکه~$R \subseteq A \times A$ ورباندې وي، کټ مټ
هغه څه دي چې د يوه دوه‌ځاييز اړوند سمبول~$\Obj P$ لرونکې د لومړۍ
درجې ژبې د !!a{structure} دپاره ئې غواړو: $\Domain{M} = A$ او
$\Assign{\Obj P}{M} = R$ ټاکو۔ ځکه~$R$ مخکښ‌مرتبه ده، انعکاسي او
انتقالي ده، او د لومړۍ درجې منطق د !!{sentence}s داسې سټ~$\Gamma$
موندلے شو چې دا خبره وايي:
\begin{align*}
  & \lforall[\Obj v_0][\Atom{\Obj P}{\Obj v_0,\Obj v_0}]\\
  & \lforall[\Obj v_0][\lforall[\Obj v_1][\lforall[\Obj v_2][((\Atom{\Obj P}{\Obj v_0,\Obj v_1} \land \Atom{\Obj P}{\Obj v_1,\Obj v_2}) \lif \Atom{\Obj P}{\Obj v_0,\Obj v_2})]]]
\end{align*}
دا !!{sentence}s يوازې د ``د هر~$x$ دپاره~$Rxx$'' ($R$ انعکاسي
ده) او ``هر کله چې~$Rxy$ او~$Ryz$ وي، نو~$Rxz$ هم وي'' ($R$
انتقالي ده) سمبولي بڼې دي۔ وينو چې !!a{structure}~$\Struct{M}$ د
دغو دوو !!{sentence}s~$\Gamma$ هله او يوازې هله مدل دے چې~$R$
(يعنې~$\Assign{\Obj P}{M}$) پر~$A$ (يعنې~$\Domain{M}$) مخکښ‌مرتبه
وي۔ په بل عبارت، د~$\Gamma$ مدلونه کټ مټ مخکښ‌مرتبې دي۔ د ټولو
مخکښ‌مرتبو هر هغه خاصيت چې د~$\Obj P$ يوازيني !!{predicate} لرونکې
د لومړۍ درجې ژبه کښې بيانېدے شي (لکه پورته انعکاسيت او انتقاليت)،
په~$\Gamma$ کښې له دوو !!{sentence}s څخه لازمېږي، او برعکس۔ نو د
~$\Gamma$ د مدلونو په باب چې څه ثابتولے شو، د ټولو مخکښ‌مرتبو په
باب مو ثابت کړي دي۔

د هرې ځانګړې تيورۍ او د مدلونو د ټولګي (لکه~$\Gamma$ او ټولې
مخکښ‌مرتبې) دپاره به په اړونده د لومړۍ درجې ژبه کښې د بيانېدونکو او
نۀ بيانېدونکو څيزونو په باب په زړه پورې پوښتنې وي۔ د
!!{structure}s ځينې داسې خاصيتونه شته چې د ټولو ژبو او د مدلونو د
ټولګيو دپاره په زړه پورې دي، يعنې د !!{domain} د کچې اړوند خاصيتونه۔
مثلاً د هر~$n \in \PosInt$ دپاره تل بيانولے شو چې !!{domain} کټ مټ
$n$~!!{element}s لري۔ د نامتناهي ډېرو !!{sentence}s د يوه سټ په
کارولو دا هم بيانولے شو چې !!{domain} نامتناهي ده۔ خو دا نۀ شو
بيانولے چې د تعريف ساحه متناهي يا !!{nonenumerable} ده۔ د لومړۍ
درجې ژبو د محدوديتونو دا پايلې د متناهي شاهد او
د لوېنهايم--سکولم له قضيو څخه راځي۔

\end{document}
